\chapter{Kernel Bypass and Low Latency Networking}
\label{chap:kernelbypass}

Chapters~\ref{chap:linuxnetworking}--\ref{chap:eventdrivenio} stay inside the 
\highlight{kernel networking stack}: sockets, skbs, interrupts, and syscalls. 
That path is correct and portable --- and often too slow or too jittery for the hottest 
market-data and order paths. \highlight{Kernel bypass} moves packet I/O into user space 
(or a thin driver path) so the trading process touches buffers with far fewer kernel 
transitions.
\newline
This chapter is conceptual.
\par\noindent\textbf{Interview tip:} focus on \highlight{why} bypass exists and 
\highlight{when} it is justified, not a full DPDK tutorial.

\section{Why Kernel Bypass}
\label{sec:why-kernel-bypass}

\subsection{Kernel Networking Overhead}
\label{subsec:kernel-networking-overhead}
A typical UDP/TCP receive on Linux roughly does:
\begin{enumerate}
    \item NIC DMA into kernel memory; interrupt or NAPI softirq runs.
    \item Protocol processing (Ethernet/IP/UDP or TCP) and socket buffer enqueue.
    \item Your thread wakes (or polls) and \texttt{recv} \highlight{copies} into a user buffer 
          (unless you use specialized zero-copy APIs).
\end{enumerate}
Each stage adds latency and \highlight{variance}: scheduling, cache misses, lock contention 
in the stack, interrupt coalescing, and noisy neighbors on the same core 
(Chapter~\ref{chap:linuxprocesses}). Throughput can still be high; 
\highlight{tail latency} is what HFT cares about.

\subsection{System Call Overhead}
\label{subsec:syscall-overhead}
See Section~\ref{sec:user-kernel-mode} for user mode vs kernel mode. 
Here, syscall overhead is framed as a networking latency cost.
\newline
Every \texttt{recv}/\texttt{send}/\texttt{epoll\char`_wait} crosses into the kernel: mode 
switch, validation, possible sleep/wake. Even a non-blocking \texttt{recv} that returns 
\texttt{EAGAIN} paid for the round trip. Busy-polling sockets 
(Chapter~\ref{chap:lowlatency}) reduces wake-ups but still uses the kernel path. 
Bypass aims to \highlight{remove} those syscalls from the packet hot path: the app polls 
NIC rings (or a user-space stack) directly.

\subsection{Context Switches}
\label{subsec:kernel-bypass-context-switches}
See Section~\ref{subsec:context-switch} for the general OS definition. 
Here, context switches are a latency cost of kernel networking and syscall-heavy I/O paths.
\newline
Interrupt handlers and softirqs can preempt your hot thread; blocking I/O forces a full 
deschedule. Isolating CPUs and pinning 
(Chapter~\ref{chap:lowlatency}, Chapter~\ref{chap:linuxprocesses}) help, but the cleanest 
fix for NIC noise is often: dedicated cores, IRQ affinity away from the strategy core, 
and a bypass stack that does not bounce through the kernel for every packet.

\section{Technologies}
\label{sec:kernel-bypass-technologies}
Vendors and open-source stacks differ in API and hardware support. The shared idea: 
\highlight{user-space packet I/O} + poll mode, often with huge pages and pinned memory.

\subsection{DPDK}
\label{subsec:dpdk}
The \highlight{Data Plane Development Kit}~\cite{dpdk-docs} is a widely cited open stack for 
poll-mode drivers (PMDs), memory pools (mbufs), and run-to-completion loops on dedicated 
cores. The application owns RX/TX rings; the kernel is largely out of the data path for 
those NICs/ports.
\newline
\textbf{Noise removed:} kernel softirq/socket path, per-packet interrupts on the hot core, 
and \texttt{recv} copies into user buffers. Ops cost: you own huge pages, NUMA pools, 
core dedication, and often the interface itself (no casual \texttt{tcpdump} / iptables story 
on a fully bound PMD port).
\par\noindent\textbf{Interview tip:} poll mode (no interrupt per packet on the hot path), huge pages, NUMA-aware 
pools (Chapter~\ref{chap:numa}), and the operational cost (you become the network stack for 
that interface --- routing, filtering, and safety are your problem).

\subsection{Solarflare OpenOnload}
\label{subsec:openonload}
\highlight{OpenOnload} (Solarflare / Xilinx / AMD lineage)~\cite{openonload-docs} accelerates 
\highlight{existing} Berkeley sockets via an LD\char`_PRELOAD user-space stack for supported 
NICs. Many apps keep \texttt{socket}/\texttt{recv}/\texttt{send} while the library short-circuits 
into user space when possible.
\newline
Appeal for trading firms: less rewrite than a full DPDK port; still hardware-specific and 
ops-heavy. Exact product names evolve.
\newline
\textbf{Noise removed:} much of the kernel stack crossing while keeping the sockets API. 
Ops cost: NIC/vendor lock-in, preload/environment discipline, and debugging a hybrid path 
when some calls still fall through to the kernel.
\par\noindent\textbf{Interview tip:} the idea is \highlight{sockets API + user-space acceleration}, 
not memorizing version strings.

\subsection{AF\_XDP}
\label{subsec:af-xdp}
\texttt{AF\char`_XDP}~\cite{af-xdp-docs} is a Linux socket family for high-performance packet 
I/O into user-space \highlight{UMEM} rings, with optional zero-copy paths and XDP programs in 
the driver/NIC. It sits between ``full kernel stack'' and ``leave Linux entirely'': you still 
integrate with the kernel's XDP/AF\char`_XDP machinery, but the hot path can avoid sk\_buff 
copies and heavy protocol processing.
\newline
\textbf{Noise removed:} sk\_buff-heavy receive and extra copies on the packet path, without 
abandoning Linux entirely. Ops cost: UMEM/ring setup, driver/XDP feature matrix, and still 
writing a packet-oriented app (not a drop-in for arbitrary TCP sessions).
\newline
Useful vocabulary: XDP for early drop/redirect in-kernel; \texttt{AF\char`_XDP} when a 
user process must consume raw/fast packet streams with lower overhead than classic sockets.

\section{Zero-Copy Market-Data Parsing}
\label{sec:zero-copy-md-parsing}
Packets on a bypass path typically land in a \highlight{user-space ring} (DPDK mbuf, 
Onload buffer, \texttt{AF\char`_XDP} UMEM slot). The descriptor owns that memory until you 
release it back to the ring. \highlight{Zero-copy} parsing means decoding the wire 
message \emph{in place} --- no \texttt{malloc}, no \texttt{memcpy} into an owned 
\texttt{std::vector} on the hot path (Section~\ref{sec:avoiding-allocations}).
\newline
\texttt{std::span} is a non-owning contiguous view (pointer + length); API details live in 
Section~\ref{subsec:container-views}. Here it is only the handle you pass into the parser 
while the ring still owns the bytes.
\newline
\textbf{Noise removed:} an extra copy (and often an allocation) between NIC buffer and 
typed fields --- plus the cache pollution of touching the payload twice. 
\textbf{Hard part:} lifetime and reclaim --- a \texttt{span} must not outlive the 
descriptor; decode, hand off what you need (or copy only the small fields you will keep), 
then return the slot.
\newline
Safer loads of possibly unaligned / endian-swapped wire bytes use \texttt{memcpy} or 
explicit byte loads into a host struct. \texttt{reinterpret\char`_cast} is shown below 
because interviews ask for it; it is correct only when alignment and layout match the 
wire, and the view's lifetime is still bound to the ring slot.
\begin{lstlisting}[language=C++]
#include <cstddef>
#include <cstdint>
#include <span>

#pragma pack(push, 1)
struct WireHeader {
    std::uint8_t  msg_type;
    std::uint16_t length;   // real feeds: convert from network endian
};
#pragma pack(pop)

// pkt views bytes still owned by the RX ring / mbuf --- do not store this pointer.
const WireHeader* header_view(std::span<const std::byte> pkt)
{
    if (pkt.size() < sizeof(WireHeader)) {
        return nullptr;
    }
    return reinterpret_cast<const WireHeader*>(pkt.data());
}

void on_rx(std::span<const std::byte> pkt)
{
    const WireHeader* h = header_view(pkt);
    if (!h) {
        return;  // truncated --- drop / count
    }
    // use h->msg_type / payload subspan; then release the ring slot
    auto body = pkt.subspan(sizeof(WireHeader));
    (void)body;
}
\end{lstlisting}
\par\noindent\textbf{Interview tip:} zero-copy is ownership discipline, not a magic type --- 
\texttt{span} views the ring; reclaim is when the view dies.

\section{When To Use Them}
\label{sec:kernel-bypass-when-to-use}
Bypass is not free: hardware lock-in, harder debugging, security/ops burden, and often a 
dedicated team. Use it when measurements say the kernel path dominates your latency budget.

\subsection{Trading}
\label{subsec:kernel-bypass-trading}
\begin{itemize}
    \item \highlightbf{Market data} --- UDP/multicast firehose where microseconds and jitter 
          matter; classic \texttt{recv} + softirqs show up in profiles.
    \item \highlightbf{Order entry} --- sometimes still kernel TCP with careful tuning 
          (\texttt{TCP\char`_NODELAY}, pinned cores); bypass when the tick-to-trade path is 
          NIC-bound and proven.
    \item Prefer proving the bottleneck (Chapter~\ref{chap:cacheperformance} profiling mindset) 
          before rewriting the I/O layer.
\end{itemize}
Decision tree (short): if profiles show kernel networking / softirq / syscall dominate 
\emph{and} the SLA is tick-to-trade, consider bypass on that NIC path. If the bottleneck is 
book logic, allocation, locks, or NUMA, fix those first --- bypass will not help.
\newline
\textbf{When \emph{not} to bypass:} control-plane / admin sockets, low message rates, teams 
without NIC/ops ownership, environments that need standard Linux tooling on the same port, 
or when tuned kernel sockets already meet the budget. Bypass is a last-mile I/O move, not a 
substitute for isolation, pinning, and allocation discipline (Chapter~\ref{chap:lowlatency}).
\par\noindent\textbf{Interview tip:} say when you would stay on kernel sockets --- most paths --- 
and when measured NIC-path latency forces DPDK / OpenOnload / \texttt{AF\char`_XDP}.

\subsection{HPC}
\label{subsec:kernel-bypass-hpc}
High-performance computing and storage fabrics use similar ideas (RDMA, user-space drivers, 
poll mode) for message rates and latency between nodes. Same trade-off: performance vs 
complexity and portability.

\subsection{Ultra-Low Latency Systems}
\label{subsec:kernel-bypass-ultra-low-latency}
When the SLA is ``kernel networking is already too slow even when tuned,'' combine:
\begin{itemize}
    \item Bypass or accelerated sockets on the NIC path,
    \item CPU isolation and busy poll (Chapter~\ref{chap:lowlatency}),
    \item Careful NUMA placement of threads and buffers (Chapter~\ref{chap:numa}),
    \item Preallocated rings and no hot-path allocations (Chapter~\ref{chap:memorymodel}).
\end{itemize}
\par\noindent\textbf{Interview tip:} kernel sockets copy and syscall; bypass polls user-space rings for lower 
median and tail latency at the cost of complexity. Name DPDK (PMD), OpenOnload-style 
socket acceleration, and \texttt{AF\char`_XDP} as the Linux-native middle ground.
